\documentclass{article}
\usepackage[utf8]{inputenc}
\usepackage[T1]{fontenc}
\usepackage{amsmath, amssymb, amsthm}
\usepackage{listings}
\usepackage{xcolor}

\title{Construções Categóricas em DLat, Bool e OrdMon}
\author{Gemini 3 Flash}
\date{Abril de 2026}

\lstset{
    language=Matlab,
    basicstyle=\small\ttfamily,
    keywordstyle=\color{blue},
    commentstyle=\color{gray},
    stringstyle=\color{purple},
    breaklines=true,
    frame=single,
    captionpos=b
}

\begin{document}

\maketitle

Este documento descreve as construções categóricas para Latices Distributivos finitos (\textbf{DLat}), Álgebras de Boolean finitas (\textbf{Bool}) e Monoides Ordenados finitos (\textbf{OrdMon}). As implementações são fornecidas em sintaxe simplificada de Octave/Matlab.

\newpage

\section{DLat: Latices Distributivos Finitos}

Os latices distributivos finitos são fechados sob todos os limites e colimites finitos.

\subsection{Construções Teóricas}
\begin{itemize}
    \item \textbf{Produtos:} O produto cartesiano $L \times M$ com operações definidas componente a componente.
    \item \textbf{Equalizadores:} O conjunto de pontos onde dois homomorfismos de latices coincidem. Como as operações são pontuais, o resultado é um sub-latice.
    \item \textbf{Coprodutos:} O latice distributivo livre sobre a união disjunta de geradores, embora para casos finitos seja frequentemente computado via categoria dual de conjuntos parcialmente ordenados (posets) finitos.
\end{itemize}

\subsection{Implementação em Octave}
\begin{lstlisting}[caption={Construções DLat}]
% Produto de dois latices
function P = dlat_product(L, M)
    % L e M sao estruturas com handles .elements e .op
    [X, Y] = meshgrid(1:length(L.elements), 1:length(M.elements));
    P.elements = [L.elements(X(:)), M.elements(Y(:))];
    P.meet = @(a, b) [L.meet(a(1), b(1)), M.meet(a(2), b(2))];
    P.join = @(a, b) [L.join(a(1), b(1)), M.join(a(2), b(2))];
end

% Equalizador de dois morfismos f e g
function E = dlat_equalizer(f, g, L)
    % Filtrar elementos onde f(x) == g(x)
    idx = arrayfun(@(i) isequal(f(L.elements{i}), g(L.elements{i})), 1:length(L.elements));
    E.elements = L.elements(idx);
    E.meet = L.meet;
    E.join = L.join;
end
\end{lstlisting}


\newpage

\section{Bool: Álgebras de Boolean Finitas}

Em \textbf{Bool}, os limites (como produtos e equalizadores) são calculados como em \textbf{DLat}, mas com o requisito adicional de preservar o complemento ($\neg$).

\subsection{Construções Teóricas}
\begin{itemize}
    \item \textbf{Produtos:} $B \times C$ usando AND, OR e NOT componente a componente.
    \item \textbf{Coequalizadores:} Calculados tomando o quociente da álgebra de Boolean pelo ideal gerado pelas diferenças dos morfismos.
\end{itemize}

\subsection{Implementação em Octave}
\begin{lstlisting}[caption={Construções Bool}]
function B_out = bool_product(B1, B2)
    % Produto cartesiano de conjuntos com operacoes Booleanas
    B_out.elements = set_product(B1.elements, B2.elements);
    B_out.not = @(a) [B1.not(a(1)), B2.not(a(2))];
    B_out.and = @(a, b) [B1.and(a(1), b(1)), B2.and(a(2), b(2))];
    B_out.or  = @(a, b) [B1.or(a(1), b(1)), B2.or(a(2), b(2))];
end
\end{lstlisting}

\newpage

\section{OrdMon: Monoides Ordenados}

Os objetos de \textbf{OrdMon} são monoides equipados com uma ordem parcial $\le$ que é compatível com a operação binária.

\subsection{Construções Teóricas}
\begin{itemize}
    \item \textbf{Produtos:} $M \times N$ com a ordem do produto: $(m, n) \le (m', n')$ sse $m \le_M m'$ e $n \le_N n'$.
    \item \textbf{Equalizadores:} O sub-monoide consistindo em elementos onde os morfismos coincidem, herdando a ordem do domínio.
\end{itemize}

\subsection{Implementação em Octave}
\begin{lstlisting}[caption={Construções OrdMon}]
function OM = ordmon_product(M1, M2)
    % Combinar monoides subjacentes
    OM.elements = set_product(M1.elements, M2.elements);
    OM.mul = @(a, b) [M1.mul(a(1), b(1)), M2.mul(a(2), b(2))];
    OM.unit = [M1.unit, M2.unit];
    % Verificacao da ordem do produto
    OM.le = @(a, b) M1.le(a(1), b(1)) && M2.le(a(2), b(2));
end

function E = ordmon_equalizer(f, g, M)
    % Filtrar elementos onde os morfismos sao iguais
    E.elements = {};
    for i = 1:length(M.elements)
        if isequal(f(M.elements{i}), g(M.elements{i}))
            E.elements{end+1} = M.elements{i};
        end
    end
    E.mul = M.mul;
    E.le = M.le;
end
\end{lstlisting}

\end{document}